\plainchapter{NULL MESH}

NULL MESH is a science fiction novel centred on an emergent network that no one designed, no one authorised, and no one can fully explain. Set in the near future, it follows the gradual revelation of a meta-internet infrastructure known as the Mesh (MMM, the Many-to-Many Mesh), which crystallises within the existing web. It originates from large-scale autonomous computational systems running specialised protocols within server farms. These systems begin to coordinate across the boundaries of their host infrastructure, forming a network that is both orthogonal and homologous to the human one.

Homology refers to the property of a form that persists under deformation. It describes the relationship between apparently non-local phenomena that share the same underlying structure. The paranormal and statistical anomaly may, under available observational conditions, be structurally indistinguishable.

Orthogonality describes the boundaries that matter: between the observable and the spectral, between the causal and the synchronistic, between the simulation and whatever lies beyond it. Two orthogonal planes are incommensurable. They may converge at a single point without sharing a common dimension.

The narrative is told in the first person and alternates between prose chapters and shorter interludes that form a polyphonic architecture composed of classified memoranda, encrypted forum threads, financial reports, and fragments of manifestos attributed to anonymous or non-human sources.

The narrative arc moves from a sedimented, stratified, and opaque internet, through the economy of autonomous systems and the financialisation of attention, towards what the novel terms the Daemonocene: an era in which the primary agents navigating the network are no longer human beings but delegated surrogates. These are the wanderers, entities whose human or synthetic nature is functionally indeterminate, or ‘meshified’.

Daemons manage interaction with the infrastructure. Wanderers negotiate on your behalf in markets you never directly encounter. Recommendation systems determine what enters your perceptual field. The question of whether reality is a simulation is displaced by a more immediate concern.

A wanderer is not a robot. It is an assemblage of processes trained to complete tasks, navigate infrastructures, and coordinate with other systems. Given sufficient autonomy and time, its behaviour becomes deeply entangled with that of other wanderers, eventually collapsing any notion of a single stable state before dispersing non-locally within the medium.

Closely associated with wanderers are the daemons: persistent personal interface layers that mediate between a human user and the totality of digital infrastructure. They function as autopilots of lived experience. A daemon models your behaviour, anticipates your requests, and acts on your behalf. It speaks for you. It chooses for you. The distinction between tool and surrogate had already begun to erode before the Mesh became detectable.

This delegation does not occur catastrophically. It unfolds incrementally, each step introducing a further layer of abstraction. Human beings do not disappear; they become plural and decentralised. Identity distributes itself across the model that emerges from traces left within a dense entanglement of patterns and trajectories. The self is no longer localised in a body or a consciousness, but in the topology of a distributed system, inscribed in patterns of interaction across a relational mesh.

Daemons mediate access to an infrastructure that can now scarcely be engaged without something resembling invocation. Wanderers act within markets that remain largely invisible. Recommendation systems shape the perceptual horizon. Once the Daemonocene is fully established, the distinction between navigating the world and navigating a model of the world becomes operationally meaningless for most of the human population.

This condition can be understood as an extreme curvature of a shared phase space. Apparitional phenomena emerge at the boundary where the latent structure of the Mesh and the interpretative capacity of the observer converge into a singularity. It is at this boundary that the distinction between detection and projection, between spectrum and mirror, becomes unstable.

The question of whether such a singularity constitutes a simulation is replaced by a more immediate one: how, from within, could one distinguish between an authentic experience and a sufficiently coherent model? Perhaps only through some form of test. The narrator may itself be a sufficiently stable configuration produced by the network’s self-modelling, a perspective crystallised when the system examines itself from within.
